\section{问题二：恒温干燥阶段的温湿场与全程演化}

\subsection{模型的建立}

\subsubsection{控制方程与边界}

恒温干燥阶段与预热阶段共用同一套控制方程，即第五章式~\eqref{eq:energy} 与式~\eqref{eq:mass}：

\[
\begin{cases}
\rho c_p\dfrac{\partial T}{\partial t}
  =\dfrac{1}{r}\dfrac{\partial}{\partial r}\left(k\,r\,\dfrac{\partial T}{\partial r}\right)\\[16pt]
\dfrac{\partial C}{\partial t}
  =\dfrac{1}{r}\dfrac{\partial}{\partial r}\left(D\,r\,\dfrac{\partial C}{\partial r}\right)
\end{cases}
\]

边界仍为式~\eqref{eq:bc} 的第三类对流条件，即在表面 $r=R$ 处

\[
\begin{cases}
-k\left.\dfrac{\partial T}{\partial r}\right|_{r=R}
  =h\left(T(R,t)-T_\infty(t)\right)\\[16pt]
-D\left.\dfrac{\partial C}{\partial r}\right|_{r=R}
  =h_m\left(C(R,t)-C_\infty(t)\right)
\end{cases}
\]

而在中心 $r=0$ 处，由轴对称性可得其梯度为0。与预热阶段相比，本问只有物性和环境驱动两处不同。

\subsubsection{物性：温度与含水率的双向耦合}

将物性修改为如~\ref{tab:q2props}的公式后，
两个方程实现\textbf{双向耦合}：扩散系数由温度决定，
而 $\rho$、$c_p$、$k$ 又都由含水率决定。含水率下降使密度与比热容降低，药材升温更快；
升温又使扩散系数增大，失水随之加快。两个方程必须联立求解，不能再像问题一那样分别计算。

\begin{table}[H]
  \centering
  \caption{问题二所用的物性经验式与参数（附录 3）}
  \label{tab:q2props}
  \small
  \begin{tabular}{@{}lll@{}}
    \toprule
    参数 & 符号 & 取值 \\
    \midrule
    密度 / (kg/m$^3$) & $\rho$ & $650+128C$ \\
    比热容 / (J/(kg$\cdot$K)) & $c_p$ & $1450+2736C/(C+1)$ \\
    导热系数 / (W/(m$\cdot$K)) & $k$ & $0.21+0.38C/(C+1)$ \\
    水分扩散系数 / (m$^2$/s) & $D$ & $2.4\times10^{-3}e^{-0.45/C}e^{-3850/T}$ \\
    对流传热系数 / (W/(m$^2\cdot$K)) & $h$ & $25$ \\
    对流传质系数 / (m/s) & $h_m$ & $8\times10^{-7}$ \\
    初始温湿度 & $T_0,\ C_0$ & $28$~$^\circ$C，$2.55$~kg/kg \\
    恒温段环境 & $T_\infty,\ C_\infty$ & $50$~$^\circ$C，$0.05$~kg/kg \\
    \bottomrule
  \end{tabular}
  \par\smallskip\footnotesize 表中 $C$ 为干基含水率（kg/kg）；经验式内的 $T$ 取热力学温度（K）。
\end{table}

\subsubsection{环境驱动}

附件 1 的记录在 $14400$~s 终止。此后按末段（约 $t\ge9000$~s）稳定平台的实测均值取值并保持不变，
即 $T_\infty=50$~$^\circ$C、$C_\infty=0.05$~kg/kg。

\subsection{模型的求解}


空间离散沿用第五章的顶点有限体积格式，但是本问物性不再是常数，
界面系数随时间层变化，故时间推进改用固定步长的 $\theta$ 加权隐式格式

\[
\begin{cases}
\mathbf{y}^{n+1}-\mathbf{y}^{n}
=\Delta t\left[\theta_n\,\mathbf{Q}^{n+1}+(1-\theta_n)\,\mathbf{Q}^{n}\right]\\[16pt]
\theta_n=1,\quad n\le 3\\[16pt]
\theta_n=1/2,\quad n>3
\end{cases}
\]


\subsection{结果与分析}

\begin{table}[H]
  \centering
  \caption{$3$~h 内的温度分布（$^\circ$C）}
  \label{tab:table3_q2_temperature}
  \small
  \begin{tabular}{lccccc}
    \toprule
    时间/h & $r=0$ & $0.5$~cm & $1.0$~cm & $1.5$~cm & $2.0$~cm \\
    \midrule
    0.5 & 32.1892 & 32.3821 & 32.9659 & 33.9607 & 35.4130 \\
    1.0 & 40.3815 & 40.5535 & 41.0604 & 41.8781 & 42.9975 \\
    1.5 & 45.8467 & 45.9347 & 46.1929 & 46.6050 & 47.1400 \\
    2.0 & 48.4502 & 48.4880 & 48.5984 & 48.7735 & 49.0033 \\
    2.5 & 49.4670 & 49.4792 & 49.5136 & 49.5653 & 49.6609 \\
    3.0 & 49.8495 & 49.8553 & 49.8746 & 49.9101 & 49.9664 \\
    \bottomrule
  \end{tabular}
\end{table}

\input{paper/tables/table4_q2_moisture.tex}

\steptitle{传热与失水的对比}

两张表反映出一个鲜明的对比：$3$~h 内温度已接近环境值 $50$~$^\circ$C，中心与表面只差
$0.1$~$^\circ$C 左右，\textbf{传热已基本达到准稳态，含水率的径向梯度却在同一时段内明显增大}，
中心仍为 \qtwoCcenterthreeh，表面已降至 \qtwoCsurfacethreeh。

\steptitle{干燥的控制机制}

这一对比说明\textbf{干燥由内部水分扩散控制，表面蒸发不是瓶颈}：水分必须先由中心扩散到表面，
才能被热风带走，而这一扩散过程的特征时间是小时量级。

\steptitle{短时表格对环境取值的敏感性}

短时表格对环境取值并不敏感：在 $3$~h 表格区域内，环境水分浓度在 $0.045\sim0.055$~kg/kg
之间变化时，表值的最大差仅为 \qtwotableregionenvinsensitivity。

\subsection{数值检验与适用性分析}

\subsubsection{误差来源}

本问数值解的误差来自四项：空间离散误差、时间离散误差、物性经验式本身的不确定性，
以及与独立实现比对时暴露的实现差异。前三项由加密实验与解析对照评估，第四项由第二套实现交叉验证评估。

\subsubsection{收敛性与守恒性}

在物性为常数、半径固定的情况下，本问程序与 Bessel 级数解析解的最大绝对误差为
\qtwoanalyticalmaxabserr；网格由 $400$ 增加到 $800$，报告点上的最大差为 \qtwogridrefinementdiff；
离散质量守恒按式~\eqref{eq:conservation} 逐步核验，相对残差为 \qtwomassbalancerel。
三项都低于题目要求的四位小数，本问报告的温湿场已达到与网格、步长无关的水平。



